\documentclass[11pt,a4paper]{article}
\usepackage{gtu-el-yazisi}
\begin{document}
\notbilgi{Nümerik Analiz}{Vize notlarım 1 -- Temel teoremler, kök bulma, interpolasyon, nümerik türev ve integral}{numerik-notlarim1-vize.pdf}

\sayfa{1}
\section{Nümerik Analiz}

\begin{teorem}[1 (Rolle Teoremi)]
  $f \in C[a,b]$ \quad ve türevlenebilir

  $f(a) = f(b) \rightarrow$ en az bir \ $x = c \in [a,b]$'de \ $f'(c) = 0$
\end{teorem}

\begin{teorem}[2 (Ortalama değer Teoremi)]
  $f \in C[a,b]$ \quad ve \quad $f$ \ $(a,b)$ diferansiyellenebilir

  $c \in (a,b) \longrightarrow f'(c) = \dfrac{f(b) - f(a)}{b - a}$
\end{teorem}

\begin{teorem}[3 (Extremum değer teoremi)]
  $f \in C[a,b]$ \qquad $x$ \okunamadi{$\in$} $c_1, c_2 \in [a,b]$

  \qquad $f(c_1) \le f(x) \le f(c_2)$ olsun

  $f$ \ $(a,b)$'de türevlenebilirse $\rightarrow$ $c_1$ ve $c_2$ noktaları aralığın uç noktaları
  yada $f'$nin sıfır olduğu noktalar.
\end{teorem}

\begin{teorem}[4 (Ara değer teoremi)]
  $f \in C[a,b]$ \quad $f(a) \cdot f(b) < 0$

  O halde öyle bir $c \in (a,b)$ bulunur ki \ $f(c) = 0$
\end{teorem}

\sayfa{2}
\subsection{Riemann manasında integrallenebilir}
\[ \int_a^b f(x)\,dx = \lim_{n \to \infty} \frac{b-a}{n} \sum_{i=1}^{n} f(x_i) \]

\subsection{Taylor polinomu ve Mclaurin}
\[
  f(x) = \underbrace{P_n(x)}_{\substack{\text{tahmini}\\ \text{değer}}}
       + \underbrace{R_n(x)}_{\substack{\text{kesim}\\ \text{hatası}}}
\]
\[ P_n(x) = \sum_{k=0}^{n} \frac{f^{(k)}(x_0)}{k!}\,(x - x_0)^k \]
\[ R_n = \frac{f^{(n+1)}(\xi)}{(n+1)!}\,(x - x_0)^{n+1} \]

$x_0 = 0$ \quad Mclaurin serisini verir

\sayfa{3}
\section{Tek değişkenli fonksiyonlar ve kök bulma yöntemi}
\subsection{1) Bisection}
\begin{align*}
  & f \in C[a,b] \quad \text{ve} \quad f(a) \cdot f(b) < 0 \\
  & a_1 = a \qquad b_1 = b \\
  & p_1 = \frac{a_1 + b_1}{2}
\end{align*}
\begin{tabular}{@{}ll@{}}
  if & $f(p_1) \cdot f(a) < 0$ \ ise \ $p \in (a, p_1)$ \\
  else if & $f(p_1)\, f(b) < 0$ \ ise \ $p \in (p_1, b)$
\end{tabular}

\medskip
$n$. step $\rightarrow [a_n, b_n] \rightarrow f(a_n) \cdot f(b_n) < 0$

\qquad $p \in (a_n, b_n)$ \qquad $p_n = \dfrac{a_n + b_n}{2}$

\underline{Hatası:} \qquad $|p_n - p| \le \dfrac{b-a}{2^n}$

\sayfa{4}
\subsection{Sabit nokta (fixed point iteration)}
\begin{align*}
  f(x) = g(x) - x &= 0 \quad \text{olsun} \\
  g(x) &= x
\end{align*}
$g : [a,b] \rightarrow \mathbb{R}$ fonksiyonu için \ $g(p) = p$ \ olacak şekilde bir $p \in [a,b]$
sayısı bulunuyorsa $p$'ye $g(x)$ fonksiyonun sabit noktası denir.
\[ p_n = g(p_{n-1}) \]

\begin{teorem}
  $g \in C[a,b]$
  \begin{enumerate}[label=\arabic*)]
    \item Eğer her $x \in [a,b]$ için $g(x) \in [a,b]$ oluyorsa $g$ fonksiyonun $[a,b]$'de en az
      bir adet sabit noktası vardır.
    \item 1 sağlanıyorsa, eğer $x \in (a,b)$ için $|g'(x)| \le k$ olacak biçimde bir
      $0 < k < 1$ bulunabiliyorsa o halde, $g$ fonksiyonun $[a,b]$ aralığında sabit noktası tektir.
    \item 1 ve 2 sağlanıyorsa, her $p_0 \in [a,b]$ başlangıç iterasyonu için $p_n = g(p_{n-1})$
      SNI $g$ fonksiyonunun $p \in [a,b]$ sabit noktasına yakınsak olur.
  \end{enumerate}
\end{teorem}

\sayfa{5}
\textcolor{kirmizikalem}{Hatası:} \ Teorem'in koşulları sağlandığında $p_n$ SNI'nun $p$ sabit noktasına
yaklaşımındaki hata
\[ |p_n - p| \le k^n \max\{p_0 - a,\ b - p_0\} \]

\subsection{Newton-Raphson}
$f : [a,b] \rightarrow \mathbb{R}$ için \ $f(x) = 0$, \ $x = p$ olsun

$f \in C^2[a,b]$ \ ve $p_0 \in [a,b]$ bir nokta

$f'(p_0) \neq 0$ \ ve \ $|p - p_0|$ küçük bir büyüklük.
\[
  0 = f(\overset{p}{x}) = f(p_0) + f'(p_0)(\overset{p}{x} - p_0)
      + \frac{f''(\xi)}{2!}(\overset{p}{x} - p_0)^2
\]
\[ p \cong p_0 - \frac{f(p_0)}{f'(p_0)} = p_1 \]
\[ \Rightarrow \quad p_n = p_{n-1} - \frac{f(p_{n-1})}{f'(p_{n-1})}, \qquad n = 1, 2, \dots \]
\[ R_n = \frac{f''(\xi)}{2!}\,(p - p_0)^2 \]

\sayfa{6}
\begin{teorem}
  $f \in C^2[a,b]$, Eğer $p \in (a,b)$ noktasında $f(p) = 0$ ve $f'(p) \neq 0$ olursa o halde öyle
  bir $\delta > 0$ vardır ki, Newton iterasyonu her $p_0 \in [p - \delta,\, p + \delta]$ başlangıç
  iterasyonu için $p$ sayısına yakınsar.
\end{teorem}

\subsection{Sekant yöntemi}
\[
  p_n = p_{n-1} - \frac{f(p_{n-1})\,(p_{n-1} - p_{n-2})}{f(p_{n-1}) - f(p_{n-2})},
  \qquad n = 2, 3, \dots
\]

\subsection{İterasyon yöntemlerinde hata analizi}
Eğer; \ $\displaystyle \lim_{n \to \infty} \frac{|p_{n+1} - p|}{|p_n - p|^{\alpha}} = \lambda$

$p_n$, $p$ sayısına $\alpha$'ıncı dereceden yakınsaktır

Sabit nokta iterasyonu lineer

Newton-raphson quadratik yakınsar (Basit kök için)

\sayfa{7}
\subsection{Lagrange interpolasyonu}
\[ P_n(x) = \sum_{k=0}^{n} L_{n,k}(x) \cdot f(x_k) \]
\[ L_{n,k} = \prod_{\substack{i=0 \\ i \neq k}}^{n} \frac{(x - x_i)}{(x_k - x_i)} \]
\[
  f(x) = P_n(x) + \underbrace{\frac{f^{(n+1)}(\xi)}{(n+1)!}\,(x - x_0)(x - x_1)\cdots(x - x_n)}_{\text{Hata}}
\]

\sayfa{8}
\subsection{Newton bölünmüş farklar interpolasyon polinomu}
\[ n_k = \prod_{i=0}^{k-1} (x - x_i) \]
ve \ $a_n = f[x_0, x_1, \dots, x_n]$ \ olmak üzere
\begin{center}
  $\displaystyle N(x) = \sum_{i=0}^{k} a_k\, n_k(x)$\footnote{Claude notu: toplam indisi muhtemelen
  $\sum_{k=0}^{n}$ olmalı.}
\end{center}

\noindent\sayfa{9}\textit{(Bu sayfa defterde yan çevrilerek yazılmış: bölünmüş farklar tablosu.)}
\begin{center}
  {\renewcommand\arraystretch{1.6}
  \begin{tabular}{cccc}
    $x_i$ & $f[x_i]$ & $f[x_i, x_{i+1}]$ & $f[x_i, x_{i+1}, x_{i+2}]$ \\
    $x_0$ & $f[x_0]$ & & \\
      & & $\boxed{\dfrac{f[x_1] - f[x_0]}{x_1 - x_0}}^{\,f[x_0, x_1]}$ & \\
    $x_1$ & $f[x_1]$ & & $\dfrac{f[x_1, x_2] - f[x_0, x_1]}{x_2 - x_0}$ \\
      & & $\dfrac{f[x_2] - f[x_1]}{x_2 - x_1}$ & \\
    $x_2$ & $f[x_2]$ & & \\
      & & $\dfrac{f[x_3] - f[x_2]}{x_3 - x_2}$ & \\
    $x_3$ & $f[x_3]$ & & \\
      & & $\dfrac{f[x_4] - f[x_3]}{x_4 - x_3}$ & \\
    $x_4$ & $f[x_4]$ & & \\
      & & $\dfrac{f[x_5] - f[x_4]}{x_5 - x_4}$ & \\
    $x_5$ & $f[x_5]$ & &
  \end{tabular}}
\end{center}

\sayfa{10}
\subsection{Hermite Polinomu}
\[
  \begin{array}{lr}
    (*)\ \ H(x_k) \Rightarrow f(x_0)\ f(x_1)\ f(x_2) \cdots f(x_n) & n+1 \\
    (**)\ H'(x_k) \Rightarrow f'(x_0)\ f'(x_1)\ f'(x_2) \cdots f'(x_n) & +\ \ n+1 \\ \cline{2-2}
    & 2n+2 \text{ bilgi}
  \end{array}
\]
\[ H_{2n+1}(x) = a_0 + a_1 x + \dots + a_{2n+1} x^{2n+1} \]
\[
  H_{2n+1}(x) = \sum_{k=0}^{n} h_{n,k}(x)\, f(x_k) + \sum_{k=0}^{n} \hat{h}_{n,k}(x)\, f'(x_k)
\]
\[ H_{n,k}(x) = \bigl[1 - 2(x - x_k)\, L'_{n,k}(x_k)\bigr]\, L^2_{n,k}(x) \]
\[ \hat{h}_{n,k}(x) = (x - x_k)\, L^2_{n,k}(x) \]
\[
  f(x) = H_{2n+1}(x) + \frac{f^{(2n+2)}(\xi)}{(2n+2)!}\,(x - x_0)^2 (x - x_1)^2 \cdots (x - x_n)^2
\]

\subsection{Kübik spline interpolasyonu}
\noindent\sayfa{11}$[a,b]$ aralığında $f(x)$ fonksiyonu var olsun
\[ a = x_0 < x_1 < \dots < x_n = b \]
\[ f(x_k), \quad k = 0, 1, \dots, n \]
\[
  S(x)
  \begin{cases}
    S_0(x) & ;\ a = x_0 \le x \le x_1 \\
    S_1(x) & ,\ x_1 \le x \le x_2 \\
    \ \vdots & \\
    S_{n-1}(x) & \phantom{,}\ x_{n-1} \le x \le x_n = b
  \end{cases}
\]
\[ S_k(x) = a_k + b_k(x - x_k) + c_k(x - x_k)^2 + d_k(x - x_k)^3 \]
\begin{enumerate}[label=\alph*)]
  \item $S_0(x_0) = f(x_0)$, \ $S_k(x_{k+1}) = f(x_{k+1}) = S_{k+1}(x_{k+1})$, \\
    $k = 0, 1, \dots, n-2$, \ $S_{n-1}(x_n) = f(x_n)$
  \item $S'_k(x_{k+1}) = S'_{k+1}(x_{k+1})$, \quad $k = 0, 1, \dots, n-2$
  \item $S''_k(x_{k+1}) = S''_{k+1}(x_{k+1})$, \quad $k = 0, 1, \dots, n-2$
  \item Şunlardan biri sağlansın
    \begin{itemize}[label={}]
      \item (doğal) {\small natural} \quad 1) \ $S''(x_0) = 0 = S''(x_n) \Longleftrightarrow
        S''_0(a) = S''_{n-1}(b) = 0$
      \item bağlanmış {\small clamped} \quad 2) \ $S'(x_0) = f'(x_0)$, \ $S'(x_n) = f'(x_n)$ \\
        \hspace*{3cm} $\Longleftrightarrow S'_0(a) = f'(a)$, \ $S'_{n-1}(b) = f'(b)$
    \end{itemize}
\end{enumerate}

\sayfa{12}
\section{Nümerik Diferansiyelleme}
\subsection{İki nokta formülü:}
\[ f'(x) = \frac{f(x_0 + h) - f(x_0)}{h} - h\,\frac{f''(\xi)}{2} \]
\[ \xi \ (x_0,\ x_0 + h) \]
geri fark tüm ileri fark formüllerinde \ $h \rightarrow -h$ \ konularak bulunur

\subsection{3 nokta uç nokta formülü}
\[ f'(x_0) = P_2'(x_0) + \frac{f'''(\xi)}{6}\,2h^2 \]
\[ P_2'(x_0) = \frac{1}{2h}\bigl[-3f(x_0) + 4f(x_0 + h) - f(x_0 + 2h)\bigr] \]

\subsection{3 nokta orta nokta formülü}
\[ f'(x_0) = \frac{1}{2h}\bigl[f(x_0 + h) - f(x_0 - h)\bigr] - \frac{h^2}{6}\,f'''(\xi) \]

\sayfa{13}
\subsection{İkinci dereceden türev}
\[ f''(x_0) = \frac{f(x_0 - h) - 2f(x_0) + f(x_0 + h)}{h^2} - \frac{h^2}{12}\,f^{(4)}(\xi) \]

\section{Nümerik integralleme}
\subsection{Trapezoid (yamuk) kuralı}
\begin{center}
\begin{tikzpicture}[line width=0.7pt, murekkep]
  \draw (0,0) -- (4,0);
  \draw (0,0.12) -- (0,-0.12);
  \draw (4,0.12) -- (4,-0.12);
  \node[below] at (0,-0.12) {$x_0 = a$};
  \node[below] at (4,-0.12) {$x_1 = b$};
  \draw (0.1,0.3) .. controls (0.6,0.55) and (1.6,0.35) .. (2,0.6)
        .. controls (2.4,0.35) and (3.4,0.55) .. (3.9,0.3);
  \node[above] at (2,0.6) {$h$};
\end{tikzpicture}
\end{center}
\[ \int_a^b f(x)\,dx = \frac{h}{2}\bigl[f(x_0) + f(x_1)\bigr] - \frac{h^3}{12}\,f''(\xi_1) \]
\[ f(x) \cong P_1(x) = \frac{x - x_1}{x_0 - x_1}\,f(x_0) + \frac{x - x_0}{x_1 - x_0}\,f(x_1) \]
\[ E_1(x) = \frac{f''(\xi(x))}{2}\,(x - x_0)(x - x_1) \qquad \xi \in (x_0, x_1) \]

\sayfa{14}
\subsection{Simpson kuralı}
\[
  \int_a^b f(x)\,dx = \frac{h}{3}\bigl[f(x_0) + 4f(x_1) + f(x_2)\bigr] - \frac{h^5}{90}\,f^{(4)}(\xi)
\]

\subsection{Kapalı Newton-Cotes formülü}
$[a,b]$ kapalı aralık
\[ x_0 = a, \quad x_n = b \]
\[ x_k = x_0 + kh \qquad h = \frac{b-a}{n} \qquad k = 0, 1, \dots, n \]
\begin{center}
  $\displaystyle P_n(x) = \sum_{k=1}^{n} f(x_k) \prod_{\substack{i=0 \\ i \neq k}}^{n} \frac{(x - x_i)}{(x_k - x_i)}$%
  \footnote{Claude notu: toplam muhtemelen $k = 0$'dan başlamalı.}
\end{center}
\[ f(x) = P_n(x) + \frac{f^{(n+1)}(\xi(x))}{(n+1)!} \prod_{k=0}^{n} (x - x_k) \]
\[ \int_a^b f(x)\,dx \cong \int_a^b P_n(x)\,dx = \sum_{k=0}^{n} \alpha_k\, f(x_k) \]
\[ \alpha_k = \int_a^b \prod_{\substack{i=0 \\ i \neq k}}^{n} \frac{x - x_i}{x_k - x_i}\,dx \]
$\displaystyle \sum_{k=0}^{n} \alpha_k\, f(x_k)$ \ $(n+1)$ adet noktadan kapalı NC formülü olsun,

\qquad $x_0 = a$, \ $x_n = b$ \ ve \ $h = \dfrac{b-a}{n}$

\sayfa{15}
Böylece $\xi \in (a,b)$ sayısı, $n$ çift ve $f \in C^{n+2}[a,b]$ olacak biçimde
\[
  \int_a^b f(x)\,dx = \sum_{k=0}^{n} \alpha_k\, f(x_k) + \frac{h^{n+3} f^{(n+2)}(\xi)}{(n+2)!}
  \cdot \int_0^n t^2 (t-1) \cdots (t-n)\,dt
\]

$n$ tek ve $f \in C^{n+1}[a,b]$
\begin{center}
  $\displaystyle \int_a^b f(x)\,dx = \sum_{k=0}^{n} \alpha_k\, f(x_k) + \frac{h^{n+1} f^{(n+1)}(\xi)}{(n+1)!}
  \cdot \int_0^n t (t-1) \cdots (t-n)\,dt$\footnote{Claude notu: muhtemelen $h^{n+2}$ olmalı.}
\end{center}

\subsection{Açık Newton-Cotes formülü {\normalsize(\okunamadi{Ara nokta})}}
$(a,b)$ açık aralık
\[ x_{-1} = a, \quad x_{n+1} = b \]
\[ x_k = x_0 + kh \qquad\qquad h = \frac{b-a}{n+2} \]
\[ k = 0, 1, \dots, n \]
\begin{center}
\begin{tikzpicture}[line width=0.7pt, murekkep]
  \draw (0,0) -- (8,0);
  \foreach \x in {0,1.2,2.6,8} \draw (\x,0.12) -- (\x,-0.12);
  \node[below] at (0,-0.12) {$x_{-1} = a$};
  \node[below] at (1.2,-0.12) {$x_0$};
  \node[below] at (2.6,-0.12) {$x_1$};
  \node[below] at (8,-0.12) {$x_{n+1} = b$};
\end{tikzpicture}
\end{center}
\[ \int_a^b f(x)\,dx = 2h\, f(x_0) + \frac{h^3}{3}\, f''(\xi_1) \]

\end{document}
